\section{绪论}\label{ch:intro}

\subsection{研究背景与意义}

【本节填写研究问题的现实背景。建议从应用场景写起：研究对象是什么，为什么重要，当前实践或已有方法遇到了什么限制。】

以农业、生态、工程或信息类论文为例，背景段可以采用“对象--需求--瓶颈--本文切入点”的顺序组织。第一段说明研究对象及其应用价值；第二段说明已有技术路线；第三段指出仍未解决的问题，并自然引出本文工作。

\subsection{国内外研究现状}

【本节填写文献综述。建议按方法类别或问题链条分小节，不要按论文逐篇罗列。每个段落以判断句开头，再用文献支撑。】

已有研究通常可以分为经验统计方法、机理模型方法和数据驱动方法三类。经验统计方法实现简单，但跨场景迁移能力有限；机理模型解释性较强，但参数反演存在不确定性；数据驱动方法能够刻画复杂非线性关系，但对训练样本和评估设计较为敏感\citep{breiman2001random,verrelst2015optical}。

\subsubsection{经验统计方法}

【三级标题示例。用于展开二级标题下的具体方向、实验步骤或子问题。三级标题会在正文中编号为 1.2.1，但目录默认只列到二级标题。】

经验统计方法通常以少量可解释变量建立目标变量的回归关系，优点是实现简单、结果直观，缺点是模型参数容易受研究区、样本范围和观测条件影响。

\subsubsection{数据驱动方法}

【三级标题示例。若某一二级标题下存在多个并列方法或实验环节，建议使用三级标题，而不是用加粗段首替代正式层级。】

数据驱动方法能够从高维特征中学习复杂非线性关系，但需要关注样本代表性、数据泄露、过拟合和跨场景验证等问题。

\subsection{研究内容与技术路线}

【本节填写本文具体做什么。建议列出 3--4 项工作，每项都能在后文章节找到对应实验或结果。】

本文围绕【研究对象】开展如下工作：

（1）构建【数据集或实验样本】，完成【预处理、配准、质量控制】等步骤；

（2）设计【模型或方法】，提取【关键特征】，建立【输入--输出】映射关系；

（3）通过【交叉验证/消融实验/对比实验】评估模型性能，分析主要误差来源；

（4）总结方法适用范围，并提出后续改进方向。

\begin{figure}[ht]
  \centering
  \begin{tikzpicture}[
    node distance=6mm,
    box/.style={draw, rounded corners=2pt, minimum width=29mm, minimum height=9mm, align=center},
    arrow/.style={-{Latex[length=2.5mm]}, thick}
  ]
    \node[box] (data) {数据获取\\与预处理};
    \node[box, right=of data] (feature) {特征构建\\与样本筛选};
    \node[box, right=of feature] (model) {模型训练\\与验证};
    \node[box, right=of model] (analysis) {结果分析\\与应用评价};
    \draw[arrow] (data) -- (feature);
    \draw[arrow] (feature) -- (model);
    \draw[arrow] (model) -- (analysis);
  \end{tikzpicture}
  \caption{本文技术路线示意图}
  \label{fig:workflow}
\end{figure}

\subsection{论文组织结构}

本文共分为五章。第一章介绍研究背景、研究现状和技术路线；第二章说明研究区、数据来源和预处理流程；第三章介绍核心方法与评价指标；第四章展示实验结果并开展讨论；第五章总结主要结论、创新点和不足。
